\documentclass[12pt]{article}
\usepackage[margin=1in]{geometry}
\usepackage{minted}
\usepackage{caption}
\usepackage{url}
\usepackage{graphicx}
\usepackage[T1]{fontenc}

\setminted{baselinestretch=0.9}

% Number supplementary floats S1, S2, ... so cross-references from the main
% text ("Supplementary Appendix Figure S1") match what is printed here.
\renewcommand{\thefigure}{S\arabic{figure}}
\renewcommand{\thetable}{S\arabic{table}}

\title{Supplementary Appendix\\
\large TreeFlow: Probabilistic Modelling and Automatic Differentiation for Phylogenetics}
\author{Swanepoel et al.}
\date{}

\begin{document}
\maketitle

\section*{Model specification}

TreeFlow provides two ways to specify phylogenetic models. For standard phylogenetic analyses, models can be defined using a YAML-based model definition format that is used as input to TreeFlow's command line interfaces. For more flexible model specification, the Python developer API allows users to directly compose TensorFlow Probability distributions into a joint distribution object.

\subsection*{YAML model definition format}

The YAML model definition format provides a concise, text-based specification of standard phylogenetic models. The format is a nested dictionary with keys at various levels corresponding to model components (tree prior, substitution model, clock model, site model), their parameters, and the prior distributions on those parameters. Figure \ref{fig:fluyaml} shows an example YAML model definition for a phylogenetic analysis of influenza sequences, specifying a coalescent tree prior, strict molecular clock, discretized Gamma site rate model, and GTR substitution model with independent Gamma priors on the relative substitution rates.

\begin{figure}[ht]
    \begin{minted}[baselinestretch=0.9,fontsize=\small]{yaml}
clock:
  strict:
    clock_rate:
      lognormal:
        loc: -2.0
        scale: 2.0
site:
  discrete_gamma:
    category_count: 4
    site_gamma_shape:
      lognormal:
        loc: 0.0
        scale: 1.0
substitution:
  gtr_rel:
    frequencies:
      dirichlet:
        concentration:
        - 2.0
        - 2.0
        - 2.0
        - 2.0
    rate_ac:
      gamma:
        concentration: 0.05
        rate: 0.05
    rate_ag:
      gamma:
        concentration: 0.05
        rate: 0.05
    rate_at:
      gamma:
        concentration: 0.05
        rate: 0.05
    rate_cg:
      gamma:
        concentration: 0.05
        rate: 0.05
    rate_gt:
      gamma:
        concentration: 0.05
        rate: 0.05
tree:
  coalescent:
    pop_size:
      lognormal:
        loc: 1.0
        scale: 1.5
    \end{minted}
    \caption{YAML model definition for the analysis of the influenza dataset. The top-level keys specify the clock model, site rate model, substitution model, and tree prior. Each model component's parameters are given prior distributions from the standard distributions available in TensorFlow Probability.}
    \label{fig:fluyaml}
\end{figure}

\subsection*{Python API model specification}

For models that go beyond the standard form supported by the YAML format, the Python developer API provides direct access to TensorFlow Probability's joint distribution framework. This allows users to compose any combination of distributions and transformations into a probabilistic graphical model. Figure \ref{fig:carnivorescode} shows an example of this flexibility: the code specifies a phylogenetic model for the carnivores dataset, with highlighted lines showing the small changes required to extend the model from a single transition-transversion ratio (kappa) to a per-lineage kappa parameter. This extension, which leverages TensorFlow's vectorization and broadcasting functionality, would be difficult to express in a fixed model definition format.

\begin{figure}[ht]
    \begin{minted}[highlightlines={19-22,41-44},baselinestretch=0.9,highlightcolor=cyan,fontsize=\small]{py}
site_category_count = 4
pattern_counts = alignment.get_weights_tensor()
subst_model = HKY()

def build_sequence_dist(tree, kappa, frequencies, site_gamma_shape):
    unrooted_tree = tree.get_unrooted_tree()
    site_rate_distribution = DiscretizedDistribution(
        category_count=site_category_count,
        distribution=Gamma(
            concentration=site_gamma_shape,
            rate=site_gamma_shape
        ),
    )
    transition_probs_tree = get_transition_probabilities_tree(
        unrooted_tree,
        subst_model,
        rate_categories=site_rate_distribution.normalised_support,
        frequencies=frequencies,
        frequencies=tf.broadcast_to(
            tf.expand_dims(frequencies, -2),
            kappa.shape + (4,)
        ),
        kappa=kappa
    )
    return SampleWeighted(
        DiscreteParameterMixture(
            site_rate_distribution,
            LeafCTMC(
                transition_probs_tree,
                expand_dims(frequencies, -2),
            ),
        ),
        sample_shape=alignment.site_count,
        weights=pattern_counts
    )

model = JointDistributionNamed(dict(
    birth_rate=LogNormal(c(1.0), c(1.5)),
    tree=lambda birth_rate: Yule(tree.taxon_count, birth_rate),
    kappa=LogNormal(c(0.0), c(2.0)),
    kappa=Sample(
        LogNormal(c(0.0), c(2.0)),
        tree.branch_lengths.shape
    ),
    site_gamma_shape=LogNormal(c(0.0), c(1.0)),
    frequencies=Dirichlet(c([2.0, 2.0, 2.0, 2.0])),
    sequences=build_sequence_dist
))
    \end{minted}
    \caption{Code to specify models for the carnivores analysis using TreeFlow's Python API. Highlighted lines show changes (from previous line) to add kappa variation across lineages. The \texttt{Sample} distribution is used to make the kappa prior an independent identically-distributed sample for each branch. The frequencies parameter needs to be manually broadcasted over the added kappa dimension.}
    \label{fig:carnivorescode}
\end{figure}

For further examples and a detailed tutorial on using TreeFlow's Python API and command line interfaces, see the TreeFlow documentation at \url{https://treeflow.readthedocs.io}.

\section*{Choice of variational approximation family}

TreeFlow provides three variational approximation families for fixed-topology models. The \emph{mean field} family treats every coordinate of the (unconstrained) parameter vector as independent. The \emph{full rank} family places a full covariance over the whole parameter vector, including every node height. The \emph{root full rank} family, used for the analyses in the main text, places a full covariance over the substitution, clock and tree-prior parameters together with the tree's root height, and treats the remaining node-height ratios as independent.

Figures \ref{fig:carnivorescomparison} and \ref{fig:h3n2comparison} show the fitted posterior marginals for one run of each family on the carnivores and influenza datasets respectively, against the same BEAST 2 reference posterior. Tables \ref{tab:carnivorescomparison} and \ref{tab:h3n2comparison} give the corresponding evidence lower bound and wall-clock runtime. ELBO values are comparable between approximations fitted to the same dataset, but not between datasets: the log likelihood is a sum over alignment patterns, of which the carnivores alignment has substantially more than the influenza alignment despite having far fewer taxa.

An approximation which achieves a higher ELBO should be an indication of a better fit to the posterior, but in practice the quality of the marginal posterior estimates may differ between approximations with similar ELBOs. For instance, on the carnivores dataset, the mean field and root full rank approximations have similar ELBOs, but the mean field approximation produces a posterior on the tree height that is substantially different from the BEAST 2 reference. The root full rank approximation, which captures correlations between the tree height and other parameters, produces a posterior that is much closer to the reference.

\begin{figure}[ht]
    \centering
    \includegraphics[width=\textwidth]{manuscript/figures/carnivores-approximation-comparison.png}
    \caption{Posterior marginals for the carnivores dataset under each of TreeFlow's three variational approximation families (solid lines), against the BEAST 2 reference posterior (dashed). One variational run per family.}
    \label{fig:carnivorescomparison}
\end{figure}

\begin{figure}[ht]
    \centering
    \includegraphics[width=\textwidth]{manuscript/figures/h3n2-approximation-comparison.png}
    \caption{Posterior marginals for the influenza (H3N2) dataset under each of TreeFlow's three variational approximation families (solid lines), against the BEAST 2 reference posterior (dashed). One variational run per family.}
    \label{fig:h3n2comparison}
\end{figure}

\begin{table}[ht]
    \centering
    \begin{tabular}{lrr}
\hline
Approximation & ELBO & Runtime (minutes) \\
\hline
Mean field & -193,601 & 14.4 \\
Full rank & -193,611 & 13.7 \\
Root full rank & -193,600 & 13.9 \\
\hline
\end{tabular}
    \caption{Evidence lower bound and wall-clock runtime for each variational approximation family on the carnivores dataset, for the runs shown in Figure \ref{fig:carnivorescomparison}. Each run is 60,000 optimization iterations.}
    \label{tab:carnivorescomparison}
\end{table}

\begin{table}[ht]
    \centering
    \begin{tabular}{lrr}
\hline
Approximation & ELBO & Runtime (minutes) \\
\hline
Mean field & -21,561 & 59.3 \\
Full rank & -21,815 & 66.0 \\
Root full rank & -21,563 & 60.1 \\
\hline
\end{tabular}
    \caption{Evidence lower bound and wall-clock runtime for each variational approximation family on the influenza (H3N2) dataset, for the runs shown in Figure \ref{fig:h3n2comparison}. Each run is 60,000 optimization iterations.}
    \label{tab:h3n2comparison}
\end{table}

\section*{Parameters that grow with the tree}

The root full rank family confines its covariance block to parameters that are shared across the whole tree. A model may also have parameters that are replicated per branch, whose dimension grows with the tree in the same way the node heights do. Examples are an uncorrelated relaxed clock and the per-lineage transition-transversion ratio of the carnivores model in Figure \ref{fig:carnivorescode}. Including these in the covariance block would restore the $O(n^2)$ cost the family exists to avoid, so TreeFlow's approximation constructor takes a \texttt{mean\_field\_vars} argument naming the variables to exclude from the block; they are then approximated independently, alongside the node-height ratios. Any posterior correlation between such a parameter and the rest of the model, for instance between a branch's rate and the heights of the nodes it separates, is therefore not captured by this family, and remains a target for future work.

\end{document}
